\section{Changing the predual}\label{sec:predual}

We now assume CH. The preceding construction produced a dual space
$E=X_0^*$ in which every infinite-dimensional subspace is distinguished by
uncountably many path functionals. Its coordinate predual $X_0$ does not
itself give the required counterexample: on each separable subspace of $E$,
the coordinate evaluations form a separable set. We construct a different
predual whose evaluations approximate all the path functionals modulo
coordinate evaluations. The evaluation criterion then excludes every
separable infinite-dimensional quotient.

\subsection{A small perturbation of a quotient map}

We first isolate the functional-analytic step. For a Banach space $E$,
write $J_E:E\to E^{**}$ for its canonical embedding.

\begin{lemma}\label{twist:lemma}
Let $Y$ be a closed subspace of $E^*$ such that the evaluation map
$E\to Y^*$ is an isometric isomorphism. Set $W=E^*/Y$ and let
$Q:E^*\to W$ be the quotient map. Suppose that a bounded linear map
$B:E^*\to W$ satisfies
\[
 \|B\|\le\tfrac1{10},\qquad B^*(W^*)\subset J_E(E),\qquad
 B(E^*)\subset V:=\overline{B(Y)}.
\]
Then $X=\ker(Q-B)$ is a Banach space, restriction induces an isomorphism
$E\to X^*$, and $\overline{B(X)}=V$.
\end{lemma}

\begin{proof}
Put $L=Q-B$ and $b=\|B\|$. Every $w\in W$ has a $Q$-lift of norm at most
$(11/10)\|w\|$. Starting with such a lift $y_0$ of $w$, lift the error
$By_0$, then the next error, and so on. The resulting series converges
because $(11/10)b<1$. Its sum is an $L$-lift of $w$ of norm at most
\begin{equation}\label{eq:lift-constant}
 C\|w\|,\qquad C=\frac{11/10}{1-(11/10)b}.
\end{equation}
In particular $L$ is onto and $Cb<1$.

Restriction to $Y$ gives the decomposition
\[
 E^{**}=Y^\perp\oplus J_E(E).
\]
Indeed the restriction of $z^{**}$ belongs to $Y^*=E$, and subtracting
its canonical image leaves an element of $Y^\perp$. Moreover $Q^*$ is
an isometric isomorphism from $W^*$ onto $Y^\perp$. Since $L$ is onto,
\[
 X^\perp=L^*(W^*)=\{Q^*w^*-B^*w^*:w^*\in W^*\}.
\]
For completeness, the first equality follows by factoring a bounded
functional vanishing on $\ker L$ through $L$; the lift bound makes the
factor bounded. The displayed subspace is a graph over $Y^\perp$ with
values in $J_E(E)$. Consequently $J_E(E)$ maps bijectively onto
$E^{**}/X^\perp=X^*$. This bijection is bounded with bounded inverse,
by the open mapping theorem. Thus restriction identifies $E$ with $X^*$.

Let $G=\overline{B(X)}\subset V$. Define
\[
 A:W\longrightarrow V/G,\qquad A(Ly)=By+G.
\]
This is well-defined because two $L$-lifts differ by an element of $X$.
The lift bound gives $\|A\|\le Cb<1$. For $y\in Y$ one has
$Ly=-By$, and hence $A(By)=-By+G$. By density, $A|_V$ is minus the
canonical quotient map $V\to V/G$. If $G\ne V$, that quotient map has
norm one, contradicting $\|A\|<1$. Therefore $G=V$.
\end{proof}

\subsection{Coding path cosets by coordinates}

Return to the space of Section~\ref{sec:norming}, with $Y=X_0$.
Let $t$ range over the countably many pure families $k\ge0$ and the mixed
family, denoted by $*$. Write $q_t$ for its Type II exponent and
$\mathcal P_t$ for its set of actual infinite paths, including their
weights and metadata. Regard a path also as its bounded functional on $E$.

\begin{lemma}\label{twist:path-map}
For each family $t$ there is a contraction
\[
 S_t:\ell_{q_t}(\mathcal P_t)\longrightarrow W,
 \qquad S_te_P=QP.
\]
\end{lemma}

\begin{proof}
Consider finitely many distinct paths. By Lemma~\ref{path:rigidity},
removing a finite prefix from each leaves pairwise disjoint weight sets.
A removed prefix has finite support and belongs to $X_0$, so it does not
change the path coset. Finite truncations of the remaining tails are
allowed crops of the original path prefixes, with their uncropped metadata
unchanged. The Type II rule bounds every rational
$\ell_{q_t}$-unit combination of them by one in $E^*$. Their pointwise
limits on $E$ give bounded full-tail representatives in
$E^*$. Applying $Q$ to those representatives, and then approximating real
coefficients, gives
\[
 \Big\|\sum_P a_P QP\Big\|\le
 \Big(\sum_P|a_P|^{q_t}\Big)^{1/q_t}
\]
for every finitely supported real family $(a_P)$. Completion gives $S_t$.
\end{proof}

Each path is specified by countably many finite ordinal supports, rational
arrays, weights and integer metadata. Under CH,
\[
 |\mathcal P_t|\le(\aleph_1)^{\aleph_0}=\aleph_1.
\]
For pure family $k$, assign distinct labels to coordinates in $C_{k+1}$;
for the mixed family assign them to coordinates in $C_1$. The latter
color is shared with pure family $0$, so partition it into two sets of
cardinality $\aleph_1$ first. This gives injections
\[
 \iota_k:\mathcal P_k\to C_{k+1},\qquad
 \iota_*:\mathcal P_*\to C_1
\]
whose ranges are mutually disjoint.

Write $s_{\rm src}(t)=s_{k+1}$ for pure family $k$, and $s_{\rm src}(*)=s_1$.
The choice of parameters gives $s_{\rm src}(t)'\le q_t$. For $y\in E^*$ define
\[
 b_t(y)=\big(y(e_{\iota_t(P)})\big)_{P\in\mathcal P_t}.
\]
The pure-coordinate bound in Lemma~\ref{ch:coordinates}, followed by
finite-dimensional duality and then a supremum over finite subsets, gives
\[
 \|b_t(y)\|_{s_{\rm src}(t)'}\le\|y\|,
 \qquad \|b_t(y)\|_{q_t}\le\|y\|.
\]
Choose positive numbers $\tau_t$ with $\sum_t\tau_t\le1/10$, and put
\[
 By=\sum_t\tau_t S_tb_t(y),\qquad y\in E^*.
\]
This series converges in $W$ and defines an operator of norm at most
$1/10$.

\begin{lemma}\label{twist:properties}
The operator $B$ satisfies
\[
 \overline{B(X_0)}=\overline{\operatorname{span}}
       \{QP:P\in\mathcal P_t,\ t\text{ a family}\}=:V,
 \qquad B(E^*)\subset V,
 \qquad B^*(W^*)\subset J_E(E).
\]
\end{lemma}

\begin{proof}
The first two claims follow from finite approximation in each
$\ell_{q_t}(\mathcal P_t)$ and the summability of the $\tau_t$.
For the reverse inclusion in the first claim, use the assigned coordinate
functional: disjointness of the label ranges gives
$B e_{\iota_t(P)}^*=\tau_t QP$.

Fix $w^*\in W^*$. The coefficient family $S_t^*w^*$ belongs to
$\ell_{q_t'}(\mathcal P_t)$ with norm at most $\|w^*\|$.
Since $q_t'\le s_{\rm src}(t)$, it also belongs to $\ell_{s_{\rm src}(t)}$ with no larger
norm. The pure-coordinate bound therefore makes
\[
 v_t=\sum_{P\in\mathcal P_t}(S_t^*w^*)(P)e_{\iota_t(P)}
\]
a norm-convergent vector series in $E$, of norm at most $\|w^*\|$.
Here, as usual, an $\ell_p$ family with finite $p$ has countable support.
The vector $v=\sum_t\tau_t v_t$ converges in $E$. Testing the definition
of $B$ against $w^*$, first on finite sums and then by continuity, gives
$B^*w^*=J_Ev$.
\end{proof}

\subsection{The obstruction to all separable quotients}

Set $X=\ker(Q-B)$. Proposition~\ref{ch:predual} and Lemmas~\ref{twist:lemma}
and~\ref{twist:properties} identify $X^*$ with $E$, and show that
$\overline{B(X)}=V$. Write $R:E\to X^*$ for this restriction isomorphism. In particular
$X$ is infinite dimensional.
The following argument explains why density of these path cosets is the
property needed for the quotient problem.

\begin{figure}[ht]
\centering
\begin{tikzpicture}[>=Stealth,node distance=13mm and 36mm]
\node (x) {$X=\ker(Q-B)$};
\node[right=of x] (v) {$V=\overline{B(X)}$};
\node[above=of x] (e) {$E^*$};
\node[above=of v] (w) {$W=E^*/X_0$};
\draw[->] (e) -- node[above] {$Q$} (w);
\draw[->] (x) -- node[below] {$B|_X$} (v);
\draw[->] (x) -- (e);
\draw[->] (v) -- (w);
\end{tikzpicture}
\caption{The new predual $X$. The square commutes because $Qx=Bx$ on
$X$, and $B(X)$ is dense in the closed span $V$ of the path cosets.
The vertical arrows are inclusions; restriction identifies $X^*$ with $E$.}
\label{fig:predual}
\end{figure}
\FloatBarrier


\begin{theorem}\label{twist:counterexample}
Under CH there is an infinite-dimensional real Banach space with no
separable infinite-dimensional quotient.
\end{theorem}

\begin{proof}
Use the space $X$ just constructed. Let $F\subset E$ be any
infinite-dimensional separable closed subspace, and put
\[
 C_F=\overline{\{y|_F:y\in X_0\}}\subset F^*.
\]
This space is separable. Indeed a countable dense subset of $F$ is
supported on a common countable set of coordinates; all other coordinate
functionals vanish on $F$. Since $X_0$ is the closed coordinate span,
its restrictions belong to the closed span of a countable family.

Restriction induces a bounded map
\[
 T_F:W\longrightarrow F^*/C_F,\qquad
 T_F(Qy)=y|_F+C_F.
\]
We claim that $T_F(V)$ is nonseparable. Otherwise its closure $M$ would
be separable. The inverse image of $M$ under the quotient map
$F^*\to F^*/C_F$ is then separable: its kernel $C_F$ is separable, and the restriction of the quotient map
to this inverse image still has the quotient norm.  More explicitly, lift
a countable dense set $(m_n)$ in $M$ to vectors $(u_n)$ in $F^*$, and
choose a countable dense set $(c_j)$ in $C_F$.  For any vector $z$ in the
inverse image, first make its quotient close to $m_n$.  The quotient norm
then makes $z-u_n$ close to a vector of $C_F$, which can in turn be
approximated by some $c_j$.  Thus the countable set
$\{u_n+c_j:n,j\ge1\}$ is dense in the inverse image.
But this inverse image contains every path restriction $P|_F$.
Theorem~\ref{branch:separation} supplies an uncountable norm-separated
family among those restrictions, a contradiction.

Because $\overline{B(X)}=V$, continuity gives
$\overline{T_F(BX)}=\overline{T_F(V)}$; hence $T_F(BX)$ is nonseparable.
For $x\in X$ we have $Qx=Bx$, so the evaluation image
$\{x|_F:x\in X\}\subset F^*$ has a nonseparable image in $F^*/C_F$.
It is itself nonseparable. Given any infinite-dimensional separable closed $F'\subset X^*$,
apply this conclusion to $F=R^{-1}(F')$.  Composition with the bounded
isomorphism $R|_F:F\to F'$ identifies the two evaluation images and
preserves norm separability.  Thus the evaluations on every such $F'$
are nonseparable. The necessary direction of Proposition~\ref{ev:criterion}
therefore rules out a separable infinite-dimensional quotient of $X$.
\end{proof}

The scalar-field transfer of Section~\ref{sec:fields} supplies a
complex counterexample as well. This is a consequence of the finite-sum
lemma there, rather than of the positive real-to-complex implication.

\begin{corollary}\label{twist:density}
Under CH, the real counterexample $X$ constructed above satisfies
$\dens X=\dens X^*=\aleph_1$. Its complexification is a counterexample
over $\C$ of norm density $\aleph_1$.
\end{corollary}

\begin{proof}
Lemma~\ref{ch:norm} makes the finitely supported rational vectors a
dense subset of $E$ of cardinality $\aleph_1$, while the coordinate
vectors $e_\alpha$, $\alpha<\omega_1$, are pairwise at distance at
least one. Hence $\dens E=\aleph_1$, and the isomorphism $E\simeq X^*$
from Lemma~\ref{twist:lemma} gives $\dens X^*=\aleph_1$.

Choose a norm-dense set $D\subset S_{X^*}$ of cardinality $\aleph_1$,
and, for each $f\in D$, choose $x_f\in S_X$ with $|f(x_f)|>1/2$.
The vectors $x_f$ have dense linear span: otherwise Hahn--Banach
supplies $g\in S_{X^*}$ vanishing on their closed span, and choosing
$f\in D$ with $\|f-g\|<1/2$ contradicts $|f(x_f)|>1/2$.
Rational linear combinations of these vectors give $\dens X\le\aleph_1$.
The space $X$ cannot be separable, since it is infinite dimensional
and would then be its own forbidden quotient. Thus $\dens X=\aleph_1$.

The complexification in Corollary~\ref{ev:fields} has underlying real
space isomorphic to $X\oplus_\infty X$, so it has the same density.
That corollary also gives the absence of separable infinite-dimensional
complex quotients.
\end{proof}
